\chapter{Related Work}
\label{chap:related}

\section{Machine Learning Models for Prediction}

The role of machine learning models in predictive analytics in the agricultural sphere is central
when working with structured data under limited computing capabilities. In this paper, three popular
and interpretable models are examined: Logistic Regression, Decision Trees, and Random Forests. They
are suitable for smallholder situations and can be used with tabular farm statistics. These models
are considered in the context of Research Question~2, which investigates how well simple machine
learning methods can forecast post-harvest loss risk with limited data.

\subsection{Logistic Regression}

Logistic Regression is among the most basic classification algorithms applicable in machine
learning. It is necessarily applicable to binary classification problems, such as identifying when
the risk of post-harvest losses is high or low. The model uses input variables to estimate the
probability of a specific outcome, mapping predictions in the range 0--1 by use of a logistic
function. The simplicity and the ability to interpret results are among the greatest benefits of
Logistic Regression. The features have a linear contribution to the prediction, which permits a
clear understanding of how separate factors affect the result.

Logistic Regression has been applied widely in agriculture and is efficient with fewer computational
requirements. As \cite{jabed2024crop} point out, less complex models like Logistic Regression are
still very applicable in situations with small datasets where interpretability is paramount. This is
more significant among the smallholders who need understandable and distinguishable decision support
systems. However, an inconvenience of Logistic Regression is that it does not capture non-linear
relationships of variables, and hence cannot describe the complexity of more intricate
agroscenarios.

\subsection{Decision Trees}

Decision Trees offer a classification and prediction method that is rule-based and hence very
intuitive and easy to interpret. This model effectively subdivides the dataset according to feature
values, creating a tree-like structure of decisions. The internal nodes signify decision rules and
the leaf nodes signify the predicted outcomes. This framework enables users to follow the
decision-making process step by step, which proves to be important especially in real-world
agriculture.

The use of Decision Trees can be accomplished effectively to work with both numerical and categorical
data, as it is flexible enough to represent the agricultural datasets which frequently have different
types of features. According to \cite{meghraoui2024applied}, tree-based models have been extensively
used in agricultural prediction activities because they are easy to interpret and can work with
heterogeneous data. Moreover, Decision Trees are able to describe non-linear trends among variables
and therefore have better prediction ability than linear models. They are however susceptible to
overfitting, particularly in situations where the tree is overly complicated, which may lower
generalisability to new data.

\subsection{Random Forest}

Random Forest is an ensemble learning technique which constructs several Decision Trees and
aggregates their predictions to enhance overall performance. The aggregation of multiple tree outputs
additionally reduces the risk of overfitting and enhances predictive performance. The trees in the
forest are trained on random subsets of the data and features, which gives diversity in the models
and makes them more robust.

Random Forest has consistently proven to outperform individual models in agricultural prediction
work. According to \cite{jabed2024crop}, ensemble methods, especially Random Forest, can be used to
capture complicated interactions between variables while being relatively accurate. As
\cite{meghraoui2024applied} further explain, these models can process a large feature set together
with noisy data, both typical of an agricultural setting. Although, unlike single Decision Trees, a
Random Forest cannot be readily interpreted, it has the ability to conduct feature importance
analysis, which gives a picture of the most impactful variables in making predictions.

\subsection{Relevance to the Study}

These three models have been chosen based on trade-offs between the need to simplify, be
interpretable, and be able to predict. Logistic Regression is transparent and easy to use, Decision
Trees offer interpretable rule-based insights and higher accuracy, and Random Forest provides an
enhanced ensemble learning prediction. Collectively, the models allow for a holistic analysis of
post-harvest loss risk, even with a limited amount of data. Within the context of this comparative
approach, Research Question~2 lies in evaluating the trade-offs between model complexity and
predictive power in the smallholder agricultural environment.


\section{Decision Support Systems in Agribusiness}

Agribusiness decision support systems (DSS) refer to computer-based systems which help farmers and
other agricultural stakeholders in coming up with sound decisions by analyzing data and providing
actionable information. These systems are able to combine data of various sources, analyze through
analytical models, and provide recommendations that aid in planning, resource allocation, and risk
management. Regarding the case of smallholder agriculture, DSS can enhance productivity, alleviate
uncertainty, and create sustainability through evidence-based decisions.

Having DSS in agriculture is essential because they can be used to analyze raw data and give
meaningful information. The nature of agricultural decision-making is complicated because of the
interplay of economic, operational, and environmental factors. This complexity can be reduced with
the help of DSS that can detect patterns and predict outcomes and therefore facilitate timely and
effective interventions. According to \cite{duan2021multicriteria}, DSS can help achieve sustainable
agribusiness through enhancing efficiency, optimizing resource utilisation, and long-term planning.
Their multicriteria assessment shows how properly designed systems can help to improve both economic
and environmental performance in farming environments.

In spite of these advantages, current DSS are usually limited, which restricts their use by
smallholder farmers in developing countries. Many systems are technologically intricate and require
good connectivity, high-quality data, and expert knowledge to operate. This poses a hindrance to
the smallholder who might not have access to reliable internet, digital capabilities, or technical
training. Moreover, some DSS target large commercial farming and are therefore less relevant to
the conditions of small farming. \cite{duan2021multicriteria} observe that technical performance is
not the only factor to consider when evaluating the use of DSS: usability and accessibility are
often disregarded in system design.

The other major weakness is the absence of contextual adaptation. Differences in agricultural
conditions across regions are high, and DSS whose recommendations do not consider local practises,
climate conditions, and the scarcity of resources can produce impractical and irrelevant
recommendations. This enhances mistrust toward the system, granting it little usefulness to the
end users.

These challenges mean that simple, low-cost, and understandable DSS that are smallholder-specific
are urgently needed. These systems must be built on easily accessible data, employ clear-cut models,
and contain clear and practical suggestions. This is the need that this study caters to, as it will
create a machine learning-based predictive model which can assist in making decisions in the
post-harvest process. Based on simplicity and usability, the proposed approach will fill the gap
between sophisticated analytics and user intervention, which directly leads to Research Question~3.


\section{Research Gap}

Although machine learning has been extensively used in the agricultural sector, there are major gaps
in its use for smallholder agribusinesses in developing nations. The existing literature has been
largely concerned with high-accuracy types of models, including typically complicated algorithms and
large-scale datasets, which may not be readily obtainable and applicable to smallholder
settings~\cite{jabed2024crop}. This creates a detachment between the use of technology and the
reality of application.

A significant gap is the absence of cheap machine learning applications that are environmentally
adapted to smallholder farmers. Most existing systems need powerful computing power, good internet
connectivity, and specialist skills, which restrict their use in resource-limited environments. It
is further noted that existing solutions are usually set to operate at large-scale commercial
farming without much regard for the operational realities and constraints of smallholder
agrobusinesses.

Another essential disparity is the underutilisation of basic, structured data. Although deep
learning approaches have demonstrated good performance, they usually require large amounts of data,
which is not frequently accessible in developing areas~\cite{meghraoui2024applied}. Little
literature has investigated how simpler machine learning models can successfully use small, tabular
data to produce meaningful predictions.

In addition, the merging of predictive models and real actionable decisions is largely absent.
Numerous reports are devoted to prediction accuracy without converting findings into practical
suggestions for farmers. In this paper, these gaps are filled by identifying an interpretable,
low-cost machine learning methodology on limited data to aid post-harvest decision making in
smallholder settings.


\section{Conceptual Framework}

This paper follows a conceptual framework that relates input variables, machine learning processes,
and decision-oriented outputs to respond to the risk of post-harvest losses. The framework is
designed around three major components: inputs, processing, and outputs.

The \textbf{input stage} involves the environmental and operational factors that affect post-harvest
results. These are variables in terms of temperature, humidity, storage method, storage duration,
time of transportation, and handling measures. Such factors are generally known to be significant
determinants of post-harvest losses in agricultural systems.

The \textbf{processing stage} includes the use of machine learning models, namely Decision Trees,
Logistic Regression, and Random Forests. These models are used to predict relationships between
input variables and historical data on losses, in order to determine trends and provide predictive
information. They are chosen according to their capability to process structured information, their
readability, and their low computational cost.

The \textbf{output stage} generates risk estimates and recommendations for action. These outputs
aid decision making by allowing farmers to detect high-risk situations and implement proper
post-harvest management strategies that mitigate losses and enhance profitability.
